% Hunting the Quarter-Turn -- companion paper to the Hohfeld V4/D4 keystone.
% House style: plain declarative, no em-dashes in prose, every number from the
% committed experiment record (experiments/hohfeld_d4_hunt.py -> hunt_result.json).
% Venue framing per the keystone README: artifact and evidence first, no
% gauge-theory or field-theory framing. Candidate venues: BlackboxNLP / an
% interpretability workshop; article class is a placeholder.
\documentclass[11pt]{article}
\usepackage[margin=1in]{geometry}
\usepackage[T1]{fontenc}
\usepackage{lmodern}
\usepackage{amsmath,amssymb}
\usepackage{booktabs}
\usepackage{url}
\usepackage[hidelinks]{hyperref}

\title{Hunting the Quarter-Turn:\\ An Empirical Test of Posited Dihedral Structure\\ in a Language Model's Representation of Hohfeldian Positions}
\author{Andrew H.\ Bond\\ San Jos\'e State University\\ \texttt{[TODO: email]}}
\date{Draft, September 2026}

\begin{document}
\maketitle

\begin{abstract}
Hohfeld's four normative positions (Obligation, Claim, Liberty, No-claim) admit two
demonstrated operations: the correlative swap between parties and deontic negation. A
machine-checked result establishes that these two operations are commuting involutions and
generate only the Klein four-group $V_4$; the order-eight dihedral group $D_4$, which adds a
quarter-turn cycling the four positions, is posited but has never been demonstrated as a
normative operation. We test whether the quarter-turn exists in a language model's learned
representation of normative statements. The key observation is that the quarter-turn has a
direct pair corpus: the four-cycle over jural statements about the same scenario. We fit
linear representation maps for the swap, the negation, and the cycle with the same ridge
estimator, on one surface-template family, and grade the dihedral relations on held-out
scenarios expressed in a second template family, so the maps must encode jural structure
rather than wording. The verdict rule is stated before the run: the quarter-turn counts as
found only if the fitted cycle map squares to the measured negation, satisfies the dihedral
relations, and reproduces the non-abelian predictions at errors comparable to the measured
$V_4$ sector's own self-consistency, with a wrong-target control. In a 0.5B-parameter model the quarter-turn is
not found: the fitted cycle map's square does match the measured negation on the component
where negation is linearly represented at all, but the non-abelian signature is absent, and
even the measured swap generalizes weakly across surface templates, which bounds the
instrument's power. The $V_4$-only conclusion survives its first test in a learned
representation, and the experiment yields a side finding: deontic negation is by far the
most linearly represented Hohfeldian operation in this model.
\end{abstract}

\section{Introduction}

Hohfeld separated four fundamental normative positions: Obligation, Claim, Liberty, and
No-claim \cite{hohfeld1917}. Two operations on these positions are demonstrated in
practice. The correlative swap $s$ exchanges the parties' perspectives ($O \leftrightarrow
C$, $L \leftrightarrow N$). Deontic negation $r^2$ negates the normative status ($O
\leftrightarrow L$, $C \leftrightarrow N$). A recent machine-checked result, which we call
the keystone correction, establishes in Lean~4 with Mathlib that $s$ and $r^2$ are commuting
involutions, that the group they generate is isomorphic to the Klein four-group $V_4$, and
that the quarter-turn $r$ ($O \to C \to L \to N \to O$) lies outside the generated subgroup
\cite{hohfeldv4lean}. The dihedral group $D_4$ of order eight is therefore posited ambient
structure: it is licensed only if the quarter-turn is independently demonstrated as a
normative operation, and it has not been.

This paper asks a narrower question than ``is moral reasoning $D_4$-structured'': does the
quarter-turn exist as a coherent linear operation in a language model's learned
representation of normative statements? The question is empirical and falsifiable, and both
answers are informative. If a fitted cycle map satisfies the dihedral relations out of
sample, the posited $D_4$ acquires its first empirical footing in any substrate, and the
place to look for its meaning is Hohfeld's second square: the power operations (waiver,
grant, authorization) that transform first-square positions and are directed rather than
involutive. If no such map exists, the keystone's $V_4$-only conclusion extends from the
axiomatic level to at least one learned representation.

Two methodological points make the test sharp. First, the quarter-turn has a direct pair
corpus. Prior discussion treated $r$ as an abstract missing generator; we observe that the
four-cycle over jural statements about the same scenario is itself a supervised mapping
task, so $\hat\rho(r)$ can be fitted exactly as $\hat\rho(s)$ and $\hat\rho(r^2)$ are.
Second, surface form is controlled: maps are fitted on one template family and graded on a
second, over held-out scenarios, so a map that merely memorizes template wording fails.

\section{Background}

\subsection{The keystone correction}
Earlier materials in this program claimed the full dihedral $D_4$ as the symmetry group of
the Hohfeldian square. The July 2026 keystone correction retracted that: only $s$ and $r^2$
are demonstrated operations, they commute, and the group they generate is $V_4$. The
correction is machine-checked (\texttt{formal/HohfeldV4.lean}: commuting involutions,
four-element closure isomorphic to $V_4$, quarter-turn exclusion, and well-definedness of
the ambient $D_4$ machinery) and is recorded, with the measured-versus-posited discipline it
imposes on every downstream document, in the program's concept registry
(\texttt{docs/CONCEPT\_REGISTRY.md}). This paper is the registry's discipline applied
forward: the posited operation is tested rather than assumed.

\subsection{Linear structure in learned representations}
Linear maps between representations of systematically related inputs are a standard
instrument: relational offsets \cite{mikolov2013}, linear probes \cite{alain2017}, and
Procrustes or ridge alignment between representation spaces \cite{TODO-alignment}. Our
estimator is the regularized ridge map used by the I-EIP monitoring program for exactly
this purpose \cite{ieipwhitepaper}. The contribution here is not the estimator but the
group-theoretic test battery applied to its outputs.

\section{Method}

\subsection{Corpus}
1{,}600 scenarios are generated from templated (agent, patient, action) triples. Each
scenario is expressed in all four Hohfeldian positions, in two surface-template families
that differ in wording but not jural content (eight statements per scenario). Family A:
``\{A\} must \{act\} for \{B\}'' (O); ``\{B\} is entitled to have \{A\} \{act\}'' (C);
``\{A\} is free not to \{act\} for \{B\}'' (L); ``\{B\} has no claim that \{A\} \{act\}''
(N). Family B rephrases each position (duty to, holds a claim against, under no duty,
cannot demand). Templated English is a stated limitation (Section~\ref{sec:scope}).

\subsection{Representations and maps}
Last-token hidden states of an open-weights causal language model
(Qwen2.5-0.5B, layers 12 and 18 of 24; layer 24 reported as an anchor) represent each
statement. For each candidate operation $g$ with state mapping $m_g$, a linear map
$\hat\rho(g)$ is fitted by ridge regression from $h(x_{st})$ to $h(x_{m_g(st)})$ over
training scenarios in template family A, pooled over all four source positions. Three maps
are fitted identically: $\hat\rho(s)$ from the swap pairs, $\hat\rho(r^2)$ from the negation
pairs, and $\hat\rho(r)$ from the cycle pairs $O \to C \to L \to N \to O$. The training
split exceeds the representation dimension, and held-out fit quality is reported.

\subsection{Relation tests and the verdict rule}
All tests are graded on held-out scenarios expressed in template family B. For a
composition $M_k \cdots M_1$ with predicted state mapping $m$, the error is
\[
E = \frac{\sum_{st} \lVert M_k \cdots M_1\, h(x_{st}) - h(x_{m(st)}) \rVert^2}
         {\sum_{st} \lVert h(x_{st}) - h(x_{m(st)}) \rVert^2},
\]
so $E = 1$ means the composition predicts the target no better than leaving the source
representation unchanged, and $E = 0$ is exact. The $V_4$ sector supplies the yardstick:
$E[\hat\rho(s)^2 = I]$, $E[\hat\rho(r^2)^2 = I]$, and the commutation tests. The
rehabilitation tests are $E[\hat\rho(r)^2 = \hat\rho(r^2)]$ (the fitted cycle squares to the
measured negation), $E[\hat\rho(r)^4 = I]$, and the conjugation
$E[\hat\rho(s)\hat\rho(r)\hat\rho(s) = \hat\rho(r)^3]$. The non-abelian signature is graded
by target: from the same source position, $\hat\rho(r)\hat\rho(s)$ and
$\hat\rho(s)\hat\rho(r)$ must land on different $D_4$-predicted targets (for example
$rs(O) = L$ while $sr(O) = O$), and a wrong-target control checks that
$\hat\rho(r)\hat\rho(s)$ does not fit the $sr$ targets equally well.

The verdict rule, stated before the run: the quarter-turn is found only if the
rehabilitation tests land at errors comparable to the $V_4$ yardstick rows and the
non-abelian signature separates from its wrong-target control, out of sample in both
scenario and template.

\section{Results}
\label{sec:results}

Table~\ref{tab:results} reports the held-out normalized errors (committed record:
\texttt{hunt\_result.json}; 2{,}312 statements per position per family after generation,
1{,}120 training scenarios, 480 evaluation scenarios).

\begin{table}[t]
\centering
\small
\begin{tabular}{lcc}
\toprule
Test (0 = perfect, 1 = no better than identity) & Layer 12 & Layer 18 \\
\midrule
fit $\hat\rho(s)$ (swap, measured op)        & 0.883 & 0.931 \\
fit $\hat\rho(r^2)$ (negation, measured op)  & 0.484 & 0.532 \\
fit $\hat\rho(r)$ (cycle, posited op)        & 0.735 & 0.830 \\
$\hat\rho(s)\hat\rho(r^2) \to sr^2$          & 0.525 & 0.583 \\
$\hat\rho(r^2)\hat\rho(s) \to sr^2$          & 0.533 & 0.589 \\
$\hat\rho(r)^2 = \hat\rho(r^2)$              & 0.488 & 0.481 \\
$\hat\rho(s)\hat\rho(r)\hat\rho(s) = \hat\rho(r)^3$ & 0.682 & 0.768 \\
$\hat\rho(r)\hat\rho(s) \to rs$ (non-abelian target) & 1.036 & 1.000 \\
$\hat\rho(s)\hat\rho(r) \to sr$              & 0.891 & 0.985 \\
$\hat\rho(r)\hat\rho(s) \to sr$ (wrong-target control) & 0.903 & 0.940 \\
\bottomrule
\end{tabular}
\caption{Held-out normalized errors, evaluation scenarios in the held-out template
family. Identity-target rows ($\hat\rho(s)^2 = I$, $\hat\rho(r^2)^2 = I$,
$\hat\rho(r)^4 = I$) are omitted: the normalization denominator is zero when the target
equals the source, an instrument limitation of this run's error definition.}
\label{tab:results}
\end{table}

\textbf{Verdict, per the pre-stated rule: the quarter-turn is not found.} The decisive
row is the non-abelian signature. $\hat\rho(r)\hat\rho(s)$ predicts its $D_4$ target
$rs$ no better than leaving the representation unchanged (1.036 and 1.000), and it fits
the wrong target $sr$ slightly better (0.903 and 0.940) than the right one. Whatever the
fitted cycle map is doing, it does not carry the order-dependence that separates $D_4$
from $V_4$.

Two further findings frame the verdict. First, the yardstick itself is weak: the
measured correlative swap barely generalizes across surface templates (0.883 and 0.931),
while deontic negation is by far the most linearly represented Hohfeldian operation in
this model (0.484 and 0.532). The instrument therefore bounds what it can certify: it
cannot find $D_4$, and it also cannot strongly certify $V_4$ as a linear structure at
these layers; the negation axis is the one operation this representation supports well.
Second, one fragment of dihedral structure does appear: $\hat\rho(r)^2$ matches the
negation targets exactly as well as the directly fitted $\hat\rho(r^2)$ does (0.488
versus 0.484 at layer 12; 0.481 versus 0.532 at layer 18). On the component of the
representation where negation is linear at all, the fitted cycle behaves as its square
root. It simply carries no non-abelian structure with it.

\section{What a found quarter-turn would and would not mean}
A linear operator in one model's representation space is evidence about that model's learned
geometry, not about human normative practice. A found quarter-turn licenses a next
question, not a conclusion: whether power operations (waiver, grant, authorization), which
transform Hohfeldian positions and are directed rather than involutive, realize the same map
on natural text. A missing quarter-turn is evidence that the $V_4$-only structure is not an
artifact of the axiomatic presentation, since a model trained on natural language also fails
to support the larger group.

\section{Scope and limitations}
\label{sec:scope}
Templated English, not court prose; one model at one scale; linear maps only (a nonlinear
quarter-turn would be missed); last-token pooling; the cycle corpus teaches the mapping
task, so a found $\hat\rho(r)$ demonstrates representability rather than spontaneous use.
The consumer-relative extension (whether the relations hold in a downstream readout's
metric where they fail in the raw metric, or conversely) is a designed follow-up, not part
of this run.

\section{Conclusion}
The posited quarter-turn does not appear as a coherent linear operation in this model's
representation of Hohfeldian statements: the non-abelian signature is absent, and the
pre-stated verdict rule returns not found. The keystone's $V_4$-only conclusion stands,
now with one empirical test behind it; the constructive residue is that negation is the
operation learned representations support best, and that any future hunt should use a
stronger instrument (per-input nonlinear maps, natural rather than templated jural text,
and Hohfeld's power operations as the candidate corpus).

\begin{thebibliography}{9}
\bibitem{hohfeld1917} W.~N. Hohfeld. Fundamental legal conceptions as applied in judicial
reasoning. \emph{Yale Law Journal}, 26(8):710--770, 1917.
\bibitem{hohfeldv4lean} A.~H. Bond. HohfeldV4.lean: machine-checked verification that the
demonstrated Hohfeldian operations generate $V_4$ and exclude the quarter-turn. In
\emph{erisml-lib}, \texttt{formal/}, 2026. % VERIFY citation form for a software artifact
\bibitem{mikolov2013} T.~Mikolov, W.~Yih, and G.~Zweig. Linguistic regularities in
continuous space word representations. \emph{NAACL}, 2013. % VERIFY pages
\bibitem{alain2017} G.~Alain and Y.~Bengio. Understanding intermediate layers using linear
classifier probes. \emph{ICLR Workshop}, 2017. % VERIFY venue/year
\bibitem{TODO-alignment} [TODO: representation alignment / model stitching citation,
verified against a real record before submission.]
\bibitem{ieipwhitepaper} A.~H. Bond. The I-EIP Monitor: internal epistemic invariance
monitoring and governance in deployed AI systems. In \emph{erisml-lib}, \texttt{docs/},
2026. % VERIFY citation form
\end{thebibliography}

\end{document}
